\chapter{Zusammenfassung}

\section{Verschiedene Textstile}

Hier sind die verschiedenen Möglichkeiten, Text zu formatieren:

\begin{itemize}
    \item \textbf{Fetter Text} für wichtige Informationen
    \item \textit{Kursiver Text} für Betonung
    \item \underline{Unterstrichener Text} für Hervorhebung
    \item \texttt{Monospace Text} für Code
    \item \textsc{Kapitälchen} für spezielle Effekte
\end{itemize}

\subsection{Hervorhebungen}

Kombinationen sind auch möglich: \textbf{\textit{Fett und kursiv}}, oder \textbf{\texttt{Fett und Monospace}}.

\section{Absätze und Abstand}

Dies ist der erste Absatz. Er behandelt wichtige Konzepte und zeigt, wie mehrere Sätze zusammenhängen.

Dies ist ein zweiter Absatz mit einer neuen Idee. Der Abstand zwischen Absätzen ist automatisch gestaltet.

\section{Zitate und Hervorhebungen}

Ein längeres Zitat kann mit der Quote-Umgebung eingeleitet werden:

\textit{``Die Mathematik ist die Sprache, in der Gott die Welt geschrieben hat.''}
--- Galileo Galilei
